% PokéDex Manager — manual de usuario
% Compile with: latexmk -xelatex -interaction=nonstopmode manual-usuario.tex
\documentclass[11pt,a4paper]{article}

\usepackage{fontspec}
\usepackage{polyglossia}
\setdefaultlanguage{spanish}

\usepackage[
  a4paper,
  top=2.5cm, bottom=2.3cm, left=2.1cm, right=2.1cm,
  headsep=0.9cm, footskip=1.1cm
]{geometry}

\usepackage[table]{xcolor}
\usepackage{graphicx}
\usepackage{booktabs}
\usepackage{tabularx}
\usepackage{array}
\usepackage{tikz}
\usepackage[most]{tcolorbox}
\usepackage{titlesec}
\usepackage{titletoc}
\usepackage{fancyhdr}
\usepackage{enumitem}
\usepackage{microtype}
\usepackage{caption}
\usepackage{float}
\usepackage{needspace}
\usepackage[hidelinks]{hyperref}

\usetikzlibrary{calc}

% ---------------------------------------------------------------------------
% Palette. Every pair used for text was checked against WCAG AA
% (4.5:1 body, 3:1 large) with the relative-luminance formula.
% ---------------------------------------------------------------------------
\definecolor{paper}{HTML}{FFFFFF}
\definecolor{ink}{HTML}{101216}
\definecolor{coverInk}{HTML}{0E1014}
\definecolor{slab}{HTML}{1A1D23}

\definecolor{pokeRed}{HTML}{C21426}
\definecolor{pokeRedDark}{HTML}{8E0E1B}
\definecolor{pokeRedLight}{HTML}{FF7A85}
\definecolor{pokeRedTint}{HTML}{FDECEE}

\definecolor{graphite}{HTML}{2A2F38}
\definecolor{steel}{HTML}{5B6472}
\definecolor{steelDark}{HTML}{4A5260}
\definecolor{mist}{HTML}{F4F5F7}
\definecolor{mistDark}{HTML}{E7E9ED}
\definecolor{coverMuted}{HTML}{B9C0CC}
\definecolor{coverFaint}{HTML}{8B94A3}

\definecolor{gold}{HTML}{8A6100}
\definecolor{goldTint}{HTML}{FFF6DC}
\definecolor{leaf}{HTML}{1F6A3E}
\definecolor{leafTint}{HTML}{E9F5EE}
\definecolor{water}{HTML}{14568F}
\definecolor{waterTint}{HTML}{E8F0F9}
\definecolor{psychic}{HTML}{8A2570}
\definecolor{psychicTint}{HTML}{F9EAF5}

% ---------------------------------------------------------------------------
% Type
% ---------------------------------------------------------------------------
\setmainfont{Charter}[
  BoldFont={Charter Bold},
  ItalicFont={Charter Italic},
  BoldItalicFont={Charter Bold Italic},
  Numbers=Lining
]
\setsansfont{Avenir Next}[
  UprightFont={Avenir Next Regular},
  BoldFont={Avenir Next Demi Bold},
  ItalicFont={Avenir Next Italic}
]
\setmonofont{Menlo}[Scale=0.86]

\newfontfamily\displayfont{Avenir Next}[
  UprightFont={Avenir Next Medium},
  BoldFont={Avenir Next Bold}
]

\color{ink}
\linespread{1.09}

% ---------------------------------------------------------------------------
% Headings
% ---------------------------------------------------------------------------
% A pokéball drawn at the size of a word: red cap, pale body, dark band and
% button. Used as the mark on every chapter opener.
\newcommand{\pokeball}[1][0.34cm]{%
  \begin{tikzpicture}[baseline=-0.32ex]
    \begin{scope}
      \clip (0,0) circle (#1);
      \fill[pokeRed] (-#1,0) rectangle (#1,#1);
      \fill[mist]    (-#1,-#1) rectangle (#1,0);
      \fill[graphite] (-#1,-0.055cm) rectangle (#1,0.055cm);
    \end{scope}
    \draw[graphite, line width=0.8pt] (0,0) circle (#1);
    \fill[paper] (0,0) circle (0.115cm);
    \draw[graphite, line width=0.8pt] (0,0) circle (0.115cm);
  \end{tikzpicture}%
}

\titleformat{\section}[display]
  {\normalfont\sffamily}
  {%
    \pokeball[0.30cm]\hspace{0.34cm}%
    {\color{pokeRed}\displayfont\bfseries\fontsize{26}{26}\selectfont\thesection}%
  }
  {0.28cm}
  {\color{ink}\displayfont\bfseries\fontsize{23}{27}\selectfont\raggedright}
  [{\vspace{0.16cm}\color{pokeRed}\rule{3.2cm}{2.2pt}}]

\titlespacing*{\section}{0pt}{0pt}{0.85\baselineskip}

\titleformat{\subsection}
  {\normalfont\sffamily\bfseries\fontsize{13}{16}\selectfont\color{graphite}}
  {\color{pokeRed}\thesubsection}{0.6em}{}
\titlespacing*{\subsection}{0pt}{1.25\baselineskip}{0.35\baselineskip}

% ---------------------------------------------------------------------------
% Running heads and footer
% ---------------------------------------------------------------------------
\newcommand{\pagepill}{%
  \tikz[baseline=(x.base)]{%
    \node[fill=pokeRed, text=paper, rounded corners=3pt, inner xsep=6pt, inner ysep=2.5pt,
          font=\sffamily\bfseries\footnotesize] (x) {\thepage};}%
}

\pagestyle{fancy}
\fancyhf{}
\renewcommand{\headrulewidth}{0pt}
\renewcommand{\footrulewidth}{0pt}
\fancyhead[L]{\sffamily\footnotesize\color{steelDark}\leftmark}
\fancyhead[R]{\sffamily\footnotesize\color{steelDark}PokéDex Manager}
\fancyfoot[C]{\pagepill}

\fancypagestyle{plain}{\fancyhf{}\fancyfoot[C]{\pagepill}%
  \renewcommand{\headrulewidth}{0pt}}

\makeatletter
\renewcommand{\sectionmark}[1]{\markboth{\thesection\quad #1}{}}
\makeatother

% ---------------------------------------------------------------------------
% Table of contents
% ---------------------------------------------------------------------------
\titlecontents{section}
  [2.6em]
  {\vspace{0.42em}\sffamily\bfseries\color{ink}}
  {\makebox[2.6em][l]{\color{pokeRed}\thecontentslabel}}
  {}
  {\titlerule*[0.6pc]{\color{steelDark}.}\contentspage}

\titlecontents{subsection}
  [2.6em]
  {\vspace{0.12em}\sffamily\small\color{steelDark}}
  {}
  {}
  {\titlerule*[0.6pc]{\color{steelDark}.}\contentspage}

% ---------------------------------------------------------------------------
% Callouts
% ---------------------------------------------------------------------------
\newtcolorbox{pokedexbox}[1]{
  enhanced, breakable,
  colback=pokeRedTint, colframe=pokeRed,
  boxrule=0pt, leftrule=3pt,
  arc=2pt, outer arc=2pt,
  left=10pt, right=10pt, top=9pt, bottom=9pt,
  fontupper=\small\color{ink},
  title={\sffamily\bfseries\footnotesize\MakeUppercase{#1}},
  coltitle=paper, colbacktitle=pokeRed,
  attach boxed title to top left={xshift=3pt, yshift=-2pt},
  boxed title style={boxrule=0pt, arc=2pt, left=6pt, right=6pt, top=2pt, bottom=2pt},
  top=16pt
}

\newtcolorbox{notabox}[2][mist]{
  enhanced, breakable,
  colback=#1, colframe=#1,
  boxrule=0pt, arc=3pt,
  left=11pt, right=11pt, top=9pt, bottom=9pt,
  fontupper=\small\color{ink},
  before upper={\par\noindent{\sffamily\bfseries\small\color{graphite}#2}\par\vspace{3pt}}
}

\newtcolorbox{tipbox}[1]{
  enhanced, colback=leafTint, colframe=leafTint,
  boxrule=0pt, arc=3pt, left=11pt, right=11pt, top=9pt, bottom=9pt,
  fontupper=\small\color{ink},
  before upper={\par\noindent{\sffamily\bfseries\small\color{leaf}#1}\par\vspace{3pt}}
}

\newtcolorbox{avisobox}[1]{
  enhanced, colback=goldTint, colframe=goldTint,
  boxrule=0pt, arc=3pt, left=11pt, right=11pt, top=9pt, bottom=9pt,
  fontupper=\small\color{ink},
  before upper={\par\noindent{\sffamily\bfseries\small\color{gold}#1}\par\vspace{3pt}}
}

\newtcolorbox{datobox}[1]{
  enhanced, colback=waterTint, colframe=waterTint,
  boxrule=0pt, arc=3pt, left=11pt, right=11pt, top=9pt, bottom=9pt,
  fontupper=\small\color{ink},
  before upper={\par\noindent{\sffamily\bfseries\small\color{water}#1}\par\vspace{3pt}}
}

\newtcolorbox{iabox}[1]{
  enhanced, colback=psychicTint, colframe=psychicTint,
  boxrule=0pt, arc=3pt, left=11pt, right=11pt, top=9pt, bottom=9pt,
  fontupper=\small\color{ink},
  before upper={\par\noindent{\sffamily\bfseries\small\color{psychic}#1}\par\vspace{3pt}}
}

% ---------------------------------------------------------------------------
% Figures
% ---------------------------------------------------------------------------
\captionsetup{
  font={sf,footnotesize,color=steelDark},
  labelfont={sf,bf,footnotesize,color=pokeRedDark},
  labelsep=period,
  justification=raggedright,
  singlelinecheck=false,
  skip=6pt
}

\newcommand{\pantalla}[3][\textwidth]{%
  \begin{figure}[H]
    \centering
    \setlength{\fboxsep}{0pt}%
    \setlength{\fboxrule}{0.6pt}%
    {\color{mistDark}\fbox{\includegraphics[width=#1]{img/#2}}}%
    \caption{#3}
  \end{figure}%
}

% ---------------------------------------------------------------------------
% Tables
% ---------------------------------------------------------------------------
\newcommand{\thead}[1]{\sffamily\bfseries\footnotesize\color{paper}#1}
\renewcommand{\arraystretch}{1.32}

\newcommand{\cuenta}[1]{{\ttfamily\footnotesize #1}}

\setlist[itemize]{leftmargin=1.15em, itemsep=2pt, topsep=4pt}
\setlist[description]{leftmargin=0pt, labelwidth=0pt, labelindent=0pt,
  style=nextline, itemsep=6pt, topsep=5pt,
  font=\sffamily\bfseries\color{graphite}}

\hypersetup{
  pdftitle={PokéDex Manager — Manual de usuario},
  pdfauthor={PokéDex Manager},
  pdflang={es}
}

\begin{document}

% ===========================================================================
% Portada
% ===========================================================================
\begin{titlepage}
\thispagestyle{empty}
\begin{tikzpicture}[remember picture, overlay]
  \fill[coverInk] (current page.south west) rectangle (current page.north east);

  % Pokéball motif, low contrast on purpose: it is decoration, not information.
  \begin{scope}
    \clip (current page.south west) rectangle (current page.north east);
    \coordinate (ball) at ($(current page.south east)+(-2.6cm,5.6cm)$);
    \fill[pokeRed, opacity=0.20] (ball) ++(-7,0) arc (180:0:7cm) -- cycle;
    \fill[coverFaint, opacity=0.09] (ball) ++(7,0) arc (0:-180:7cm) -- cycle;
    \draw[coverFaint, opacity=0.32, line width=1.1pt] (ball) circle (7cm);
    \draw[coverFaint, opacity=0.32, line width=1.4pt] ($(ball)+(-7,0)$) -- ($(ball)+(7,0)$);
    \fill[coverInk] (ball) circle (2cm);
    \draw[coverFaint, opacity=0.45, line width=1.6pt] (ball) circle (2cm);
    \draw[coverFaint, opacity=0.45, line width=1.2pt] (ball) circle (1.05cm);

    \fill[pokeRed, opacity=0.22] ($(current page.north west)+(-2cm,1.5cm)$) circle (5.6cm);
    \draw[coverFaint, opacity=0.16, line width=0.9pt]
      ($(current page.north west)+(-2cm,1.5cm)$) circle (7.8cm);
  \end{scope}
\end{tikzpicture}

\vspace*{3.4cm}
\noindent
{\sffamily\bfseries\fontsize{10}{12}\selectfont\color{pokeRedLight}%
 \MakeUppercase{Manual de usuario}}

\vspace{0.55cm}
\noindent
{\displayfont\bfseries\fontsize{50}{54}\selectfont\color{paper}%
 Poké\textcolor{coverMuted}{Dex}\\ Manager}

\vspace{0.7cm}
\noindent
\begin{minipage}{11.2cm}
{\sffamily\fontsize{12.5}{18}\selectfont\color{coverMuted}%
 Cataloga tus cartas, mira cuánto valen de verdad, y cámbialas con quien tenga
 justo lo que te falta.}
\end{minipage}

\vfill

\noindent
{\color{coverFaint}\rule{5.5cm}{0.8pt}}

\vspace{0.5cm}
\noindent
\begin{minipage}{\textwidth}
\sffamily\footnotesize
\setlength{\tabcolsep}{0pt}
\begin{tabular}{@{}p{3.3cm}p{9.5cm}@{}}
{\color{coverMuted}Versión} & {\color{paper}1.0} \\[3pt]
{\color{coverMuted}Aplicación} & {\color{paper}PokéDex Manager (web)} \\[3pt]
{\color{coverMuted}Pruébala en} & {\color{paper}\url{https://web-production-78a4b.up.railway.app}} \\[3pt]
{\color{coverMuted}Código} & {\color{paper}\url{https://github.com/Skulltulaidm/pokedex-manager}} \\[3pt]
{\color{coverMuted}Idioma} & {\color{paper}Español} \\
\end{tabular}
\end{minipage}

\vspace{0.9cm}
\end{titlepage}

% ===========================================================================
% Índice
% ===========================================================================
\thispagestyle{plain}
\noindent
{\sffamily\bfseries\fontsize{9.5}{11}\selectfont\color{pokeRed}%
 \MakeUppercase{Contenido}}

\vspace{0.25cm}
\noindent{\color{pokeRed}\rule{3.2cm}{2.2pt}}
\vspace{0.9cm}

\setcounter{tocdepth}{1}
\makeatletter\@starttoc{toc}\makeatother

\vspace{1.2cm}
\begin{notabox}{Cómo leer este manual}
Cada capítulo se puede leer suelto. Si sólo quieres ver la aplicación
funcionando, ve al capítulo~2: hay tres cuentas de prueba ya cargadas y una
dirección donde abrirlas sin instalar nada.
\end{notabox}

\clearpage

% ===========================================================================
\section{Qué es y qué resuelve}
% ===========================================================================

Una caja de zapatos con cartas es ilegible. No sabes qué tienes, no sabes qué
vale, y no sabes a cuántas cartas estás de terminar un set. PokéDex Manager
convierte esa caja en algo que se puede consultar.

Registras tus cartas ---a mano o fotografiándolas--- y a partir de ahí la
aplicación te responde tres preguntas que antes tenías que responder de memoria:

\begin{description}
  \item[Qué tengo.] El catálogo completo con tus cartas encendidas y las que te
  faltan apagadas, set por set.
  \item[Cuánto vale.] Precio de mercado de cada carta, cuánto pagaste, y por
  tanto cuánto has ganado o perdido.
  \item[Con quién puedo cambiar.] Otros coleccionistas que tienen repetido lo
  que buscas y buscan lo que a ti te sobra.
\end{description}

Encima de todo eso hay un asistente al que le puedes preguntar en español
---«¿cuál es mi carta más cara?», «¿qué me falta del Base Set?»--- y que se
acuerda de lo que le cuentas.

\begin{pokedexbox}{Lo que no hace}
No compra ni vende cartas, no cobra comisiones y no mueve nada por su cuenta.
Un trueque se cierra entre dos personas; la aplicación sólo lo propone, lo
valora y guarda la conversación.
\end{pokedexbox}

% ===========================================================================
\clearpage
\section{Los tres accesos}
% ===========================================================================

La forma más rápida de entender la aplicación es entrar con una cuenta que ya
tenga datos dentro. Hay tres preparadas, cada una pensada para enseñarte una
cara distinta del producto. El mismo formulario de entrada crea una cuenta
nueva: sólo hay que tocar «¿Primera vez? Crear una cuenta».

\medskip
\noindent
{\sffamily\bfseries\small\color{graphite}Ábrelas sin instalar nada en:}\\[2pt]
\noindent{\ttfamily\small\href{https://web-production-78a4b.up.railway.app}{https://web-production-78a4b.up.railway.app}}

\medskip
\noindent
{\sffamily\bfseries\small\color{graphite}Contraseña de las tres:}
{\ttfamily\small pokedex2026}

\vspace{0.5cm}

\noindent
\begin{tabularx}{\textwidth}{@{}>{\raggedright\arraybackslash}p{4.35cm} >{\raggedright\arraybackslash\sffamily\small}p{2.75cm} X@{}}
\rowcolor{graphite}
\thead{Correo} & \thead{Nombre} & \thead{Qué vas a ver} \\
\cuenta{coleccionista@pokedex.test} & AshKetchum99 &
  34 cartas en 4 sets, todas con precio de compra. Es la cuenta de la cartera:
  valor, rendimiento, posiciones y concentración. \\
\rowcolor{mist}
\cuenta{trader@pokedex.test} & xX\_TradeLord\_Xx &
  18 cartas, 9 de ellas repetidas, 10 deseos y 3 publicaciones en el tablón. Es
  la cuenta del trueque: emparejamiento, ofertas, tablón y simulador. \\
\cuenta{novato@pokedex.test} & PikaRookie &
  2 cartas y 4 deseos. Es la cuenta de empezar de cero: pantallas vacías,
  agregar la primera carta, escanear. \\
\bottomrule
\end{tabularx}

\vspace{0.35cm}

\begin{tipbox}{Por dónde empezar}
Si tienes cinco minutos, entra con \cuenta{trader@pokedex.test}: es la única
desde la que funciona la propuesta de trueque de la IA, porque necesita cartas
repetidas y una lista de deseos a la vez. Si lo que te interesa es el dinero,
entra con \cuenta{coleccionista@pokedex.test}.
\end{tipbox}

\begin{avisobox}{Son cuentas compartidas}
Cualquiera que lea este manual entra con las mismas credenciales, así que lo
que hagas en ellas lo verán los demás. Para tus cartas de verdad, créate una
cuenta propia.
\end{avisobox}

% ===========================================================================
\clearpage
\section{Tu colección}
% ===========================================================================

\subsection{El catálogo}

La pantalla principal no es tu colección: es el catálogo entero. Tus cartas
salen a color y las que no tienes salen apagadas, de modo que el hueco se ve
sin tener que buscarlo. Arriba está lo que ya vale tu cartera y lo que costaría
completar el catálogo.

Puedes filtrar por \emph{Todas}, \emph{Tuyas} o \emph{Te faltan}, ordenar,
acotar por set o generación, y usar la fila de tipos para quedarte con uno solo.
Al tocar una carta se abre su ficha: precio, rareza, PS, año, estadísticas de la
especie y su línea evolutiva.

\pantalla{catalogo-desktop.png}{Catálogo con la cuenta \cuenta{coleccionista}.
Las cartas apagadas son las que faltan; la barra de arriba compara lo que vale
tu cartera con lo que cuesta completar el catálogo.}

\subsection{Agregar una carta a mano}

Desde \emph{Agregar una carta} buscas por nombre, eliges la impresión correcta
---la imagen te dice de qué set es--- y rellenas estado, cuántas tienes, cuánto
pagaste por cada una y una nota.

\begin{datobox}{El precio de compra es opcional, pero cuenta}
Si lo dejas vacío la carta entra igual en tu colección y en el valor estimado,
pero queda fuera del rendimiento. La aplicación te lo avisa en el propio
formulario.
\end{datobox}

\pantalla{agregar-formulario-desktop.png}{Alta manual de una carta. Poner
«2» en \emph{Cuántas tienes} es lo que convierte la segunda copia en una
repetida disponible para trueque.}

\subsection{Escanear}

Le haces una foto a la carta y el modelo transcribe lo que está impreso: nombre,
número, total del set y PS. Con eso la aplicación decide qué carta es y te
enseña \emph{por qué} coincide. Si no acertó, «No es esa» abre las demás
candidatas con su motivo. Tú eliges el estado y confirmas.

\pantalla{escaneo-resultado-desktop.png}{Resultado de un escaneo. Arriba, lo que
el modelo leyó en la foto; abajo, la carta propuesta y en qué coincide.}

% ===========================================================================
\clearpage
\section{Tu cartera}
% ===========================================================================

\emph{Resumen} tiene cuatro pestañas. La primera es la foto general de la
colección; la segunda la abre carta por carta.

\begin{description}
  \item[Valor estimado.] La suma del precio de mercado de tus cartas. Debajo,
  cuántas de tus cartas tienen precio: si son 34 de 34, el total es completo.
  \item[Rendimiento.] Lo que pagaste frente a lo que valen hoy, en dinero y en
  porcentaje.
  \item[Tus sets.] Cada casilla es una carta impresa del set y las encendidas
  son las tuyas. Van ordenados por lo que cuesta terminarlos, no por cuántas
  faltan.
\end{description}

\begin{avisobox}{Por qué una carta puede no aparecer en el rendimiento}
El rendimiento compara dos cifras: lo que pagaste y lo que vale. Si nunca
registraste cuánto pagaste por una carta, esa resta no se puede hacer, así que
la carta queda fuera del cálculo ---no entra como cero, que haría parecer que
la regalaron---. Sí sigue contando en el valor estimado. Por eso una colección
puede valer \$3.339 y tener un rendimiento calculado sobre menos cartas: en la
tabla de posiciones esas aparecen con un guion en la columna de costo.
\end{avisobox}

\pantalla{resumen-desktop.png}{Resumen de \cuenta{coleccionista}: valor
estimado, cobertura de precios, rendimiento y las cartas que más pesan.}

\subsection{Posiciones y concentración}

Cada carta es una posición: sus copias juntas, lo que pagaste, lo que valen hoy,
la ganancia o pérdida y qué porcentaje de la cartera representa. Se ordena
tocando cualquier columna.

Al lado, \emph{Concentración} responde a una pregunta que el total esconde: en
cuántas cartas está la mitad de tu dinero. Una cartera que se apoya en tres
cartas corre un riesgo distinto de una repartida entre cien.

\pantalla{posiciones-desktop.png}{Posiciones ordenadas por valor, y a la derecha
la concentración: la mitad de esta cartera está en 7 de sus 34 cartas.}

% ===========================================================================
\clearpage
\section{Trueques}
% ===========================================================================

\subsection{Emparejamiento}

Un trueque necesita las dos mitades: que alguien quiera una de tus repetidas y
que tenga repetida una de tu lista de deseos. Cuando las dos coinciden, el par
aparece en \emph{Coincidencias} ya armado y valorado, con el saldo a tu favor o
en tu contra.

Desde ahí puedes proponerlo tal cual, o abrir «Elegir otras cartas» y montar
otro reparto carta por carta.

\pantalla{coincidencias-desktop.png}{Coincidencias: la aplicación arma el
intercambio y calcula la diferencia de valor de los dos lados.}

\subsection{El tablón abierto}

No siempre quieres elegir con quién. Una publicación del tablón dice qué das y
qué pides, sin destinatario, y la toma quien pueda cumplirla. Se queda en pie
hasta que alguien la acepte o la retires, y ninguna carta se mueve de tu
colección mientras tanto.

\pantalla{tablon-desktop.png}{Una publicación ajena en el tablón. El botón
aparece atenuado con el motivo al lado: falta una carta repetida de las que
pide.}

\subsection{Ofertas y contraofertas}

Una oferta dirigida llega a \emph{Trueques} y a \emph{Novedades}. Tienes cuatro
respuestas: aceptar, rechazar, contraofertar o escribir. La contraoferta abre el
mismo selector de cartas y, al enviarla, rechaza la original: no pueden quedar
dos ofertas en pie sobre la misma mesa.

\pantalla{ofertas-desktop.png}{Arriba, una oferta recibida con sus cuatro
respuestas. Abajo, una oferta que tú enviaste y que sigue pendiente.}

\begin{tipbox}{Aceptar no mueve cartas}
Aceptar cierra el acuerdo, no el intercambio físico. Las cartas cambian de dueño
cuando ustedes se ven; después cada uno actualiza su colección.
\end{tipbox}

\subsection{Mensajes}

Un trueque se cierra hablando. Desde cualquier oferta, publicación o perfil hay
un botón \emph{Mensaje} que abre la conversación con esa persona.

\pantalla{mensajes-hilo-desktop.png}{Conversación entre dos coleccionistas.
El icono de arriba a la derecha lleva directo a armar un trueque con esa
persona.}

\subsection{El simulador y la propuesta de la IA}

El simulador arma un intercambio hipotético y te enseña cómo quedaría tu cartera
después: valor, cuántas cartas distintas tendrías, en cuántas quedaría la mitad
del valor y qué porcentaje pesaría la carta más cara. No envía nada.

El botón \emph{Que la IA proponga uno} lo rellena solo. Busca entre tus
repetidas y tu lista de deseos un cambio que te convenga, y explica por qué.
Necesita las dos cosas a la vez, así que sólo funciona desde una cuenta que
tenga repetidas y deseos ---la de \cuenta{trader}, por ejemplo---.

\pantalla[12cm]{simulador-desktop.png}{Propuesta generada por la IA, con su motivo escrito arriba.}

% ===========================================================================
\clearpage
\section{El asistente}
% ===========================================================================

\emph{Preguntar} es un chat en español que lee tu colección. No inventa
precios: consulta los mismos datos que ves en pantalla y cita las cartas con su
miniatura y su precio, de modo que puedes comprobar cada cifra.

\begin{iabox}{Dos preguntas reales, con su respuesta}
\textbf{«¿Cuánto vale mi colección y en qué sets está repartida?»}\\
\emph{Tu colección tiene un valor total estimado de 3339.19 dólares y cuenta con
34 cartas distribuidas en cuatro expansiones: Jungle (10 cartas valoradas en
1241.56 dólares), Fossil (9 cartas, 1198.26 dólares), Paldea Evolved (10 cartas,
453.36 dólares) y Base Set (5 cartas, 446.01 dólares).}

\medskip
\textbf{«¿Qué me falta para completar el Base Set y cuánto me costaría?»}\\
\emph{Actualmente tienes 5 de las 102 cartas de Base Set (Blastoise, Clefairy,
Gyarados, Mewtwo y Zapdos), por lo que te faltan 97 cartas para completarlo.
Adquirir las cartas restantes te costaría aproximadamente 1633.66 dólares según
los precios de mercado actuales.}
\end{iabox}

\pantalla{preguntar-desktop.png}{El asistente cita las cartas dentro de la
respuesta: cada chip lleva a la ficha de esa carta.}

\subsection{Lo que recuerda de ti}

Si le dices algo que vale para siempre ---«sólo colecciono cartas de la primera
generación», «no gasto más de 100 dólares por carta»--- lo guarda y lo aplica en
las siguientes conversaciones. Puedes ver esa lista, y borrar lo que quieras, en
\emph{Tu perfil} $\rightarrow$ \emph{Lo que sabe el asistente}.

% ===========================================================================
\clearpage
\section{Novedades y tu cuenta}
% ===========================================================================

\emph{Novedades} reúne lo que pasó sin ti: ofertas que te llegaron, trueques que
se cerraron, mensajes nuevos y cambios de precio en cartas de tu lista de
deseos. El filtro \emph{Sin responder} deja sólo lo que espera una acción tuya,
y la campana de la cabecera lleva la misma cuenta.

\pantalla{novedades-desktop.png}{Novedades con dos entradas pendientes. El punto
rojo marca lo que todavía no has respondido.}

\subsection{Tu perfil}

En \emph{Tu perfil} cambias tu nombre visible ---el que aparece en tus ofertas---
y controlas qué ven los demás. Tu perfil público muestra tu nombre, los sets que
coleccionas y tus repetidas; nunca cuánto vale tu colección. Aparte puedes crear
un enlace público y revocable con tus cartas, sin totales de dinero. El avatar
se genera a partir de tu correo: no hay nada que subir.

% ===========================================================================
\clearpage
\section{Preguntas frecuentes}
% ===========================================================================

\begin{description}
  \item[¿De dónde salen los precios?]
  Del precio de mercado de TCGplayer, leído una vez al día. Es orientativo, no
  una tasación: el estado real y la edición cambian lo que vale una carta.

  \item[Registré una carta y no aparece en el rendimiento.]
  Le falta el precio de compra. Ábrela desde \emph{Posiciones} o desde el
  catálogo y rellena «Cuánto pagaste por cada una».

  \item[¿Por qué no me sale ninguna coincidencia?]
  Porque falta una de las dos mitades. Necesitas al menos una carta repetida
  ---cantidad 2 o más--- y al menos un deseo apuntado.

  \item[¿Qué diferencia hay entre una oferta y una publicación del tablón?]
  La oferta va dirigida a una persona concreta. La publicación no tiene
  destinatario: la toma quien pueda cumplirla.

  \item[Acepté un trueque, ¿ya cambiaron las cartas de dueño?]
  No. Aceptar cierra el acuerdo; el intercambio es físico. Después de verse,
  cada uno actualiza su colección a mano.

  \item[¿El asistente ve mis notas privadas?]
  Lee tu colección para responderte a ti. Lo que sale de tu cuenta es sólo lo
  que compartes: el perfil público y el enlace, y ninguno de los dos incluye
  notas, fechas de compra ni totales de dinero.

  \item[Escaneé una carta y acertó la especie pero no la impresión.]
  Toca «No es esa» y elige entre las candidatas. Cada una dice en qué coincidió;
  el número y el total del set son las señales que más ayudan, así que vale la
  pena que salgan nítidos en la foto.

  \item[¿Puedo sacar mis datos?]
  Sí: la colección se exporta en CSV y en JSON.
\end{description}

\vfill

\begin{notabox}{Sobre este manual}
Todas las capturas son de la aplicación funcionando, tomadas con las tres
cuentas de prueba en un navegador de 1440 px de ancho. Las cifras que aparecen son las de esas
cuentas: en la versión publicada pueden variar, porque el catálogo cargado no es
el mismo.
\end{notabox}

\end{document}
